% Appendix: Extended Experimental Details
% Intended for \appendix section of the main paper.
% Include with: % Appendix: Extended Experimental Details
% Intended for \appendix section of the main paper.
% Include with: % Appendix: Extended Experimental Details
% Intended for \appendix section of the main paper.
% Include with: % Appendix: Extended Experimental Details
% Intended for \appendix section of the main paper.
% Include with: \input{appendix_experimental_details}

\section{Extended Experimental Details}
\label{app:experimental-details}

\subsection{Hardware and Software}
\label{app:hardware}

All experiments were run on a single workstation with 62\,GB DDR4 RAM and an
Intel Core i7-6850K (6 cores / 12 threads, 3.6\,GHz, AVX2, no native BF16).
Four installed GTX~1080~Ti GPUs (compute capability 6.1) are incompatible with
the installed PyTorch build (2.12.0+cu130, which requires CC~$\geq$~7.5); all
computation is therefore CPU-only.
BF16 matrix operations fall back to FP32 via oneDNN.
Model weights are kept in \texttt{bfloat16} throughout
(31\,GB resident memory for DeepSeek-V2-Lite).

\medskip\noindent
\textbf{Software versions:}
Python~3.12,
PyTorch~2.12.0,
Transformers~5.3.0,
safetensors~0.7.0,
datasets (Salesforce/wikitext).
All forward passes use \texttt{OMP\_NUM\_THREADS=12} and \texttt{use\_cache=False}.

% ----------------------------------------------------------------
\subsection{Model Loading and the \texttt{\_tied\_weights\_keys} Bug}
\label{app:loading-bug}

Loading DeepSeek-V2-Lite via \texttt{AutoModelForCausalLM.from\_pretrained}
under Transformers~5.x silently fails to load any weights from the checkpoint.
The root cause is in the model's remote code (\texttt{modeling\_deepseek.py}):

\begin{verbatim}
_tied_weights_keys = ["lm_head.weight"]   # old Transformers 4.x list format
\end{verbatim}

Transformers~5.x calls \texttt{.keys()} on this attribute, expecting a
\texttt{dict}; the \texttt{list} raises \texttt{AttributeError}.
The exception is caught internally and stored as a \texttt{SkipParameters}
conversion error, after which \texttt{\_initialize\_missing\_keys} calls
\texttt{\_init\_weights}, re-initialising every \texttt{nn.Linear} and
\texttt{nn.Embedding} weight to $\mathcal{N}(0,\,0.02)$.
\texttt{nn.LayerNorm} weights are not handled by \texttt{\_init\_weights}
and remain at the PyTorch default (all-ones, $\sigma=0$).
The result is that all 5{,}291 named parameters take initialisation values
rather than the trained checkpoint.

\medskip\noindent
\textbf{Symptom.}
WikiText-2 perplexity $\approx 129{,}000$.
A standard detection heuristic---flagging parameters whose empirical standard
deviation satisfies $|\sigma - 0.02| < 0.003$---corrected the
\texttt{nn.Linear} and \texttt{nn.Embedding} parameters but missed the 82
LayerNorm weights, whose trained standard deviations ($\approx 0.017$--$0.019$)
lie within the same tolerance window as the initialisation value.
After heuristic repair, perplexity remained $\approx 1.7 \times 10^{11}$,
indicating the LayerNorm weights are load-bearing.

\medskip\noindent
\textbf{Fix.}
After \texttt{from\_pretrained}, every named parameter is unconditionally
overwritten by re-reading directly from the safetensors shards:

\begin{verbatim}
for shard_file, keys in sorted(shard_to_keys.items()):
    with safe_open(shard_path, framework="pt", device="cpu") as f:
        for key in keys:
            model_params[key].data.copy_(
                f.get_tensor(key).to(dtype))
\end{verbatim}

\noindent
Shards are processed sequentially; peak extra memory is one open shard
($\approx$8\,GB) alongside the resident model ($\approx$31\,GB), well within
the 62\,GB budget.
Correctness was verified by comparing per-parameter standard deviations
against checkpoint values for a representative set spanning all weight types
(embeddings, attention projections, MoE experts, layer norms, LM head).
After the fix, WikiText-2 perplexity is 6.275 (Table~\ref{tab:exp3-results}).

% ----------------------------------------------------------------
\subsection{Perplexity Evaluation Protocol}
\label{app:ppl}

\textbf{Dataset.}
WikiText-2 validation split~\citep{merity2017pointer}, loaded via the
Salesforce/wikitext HuggingFace repository.
The full validation text is concatenated into a single token sequence
(${\approx}273$K tokens under the DeepSeek-V2-Lite tokenizer).

\medskip\noindent
\textbf{Sliding-window protocol.}
Context window $L = 512$, stride $s = 512$ (non-overlapping windows).
Each forward pass scores the final $\min(s, \text{chunk length})$ tokens.
No overlap is used; the estimate is therefore a strict upper bound on the
true held-out NLL but is consistent across all experimental conditions.

\medskip\noindent
\textbf{Sample budget.}
For Experiment~2 (activation alignment), forward passes cover the first
50{,}000 tokens (${\approx}97$ chunks).
For Experiment~3 (clamping), each PPL estimate uses the first 30 chunks
(${\approx}15{,}360$ tokens, ${\approx}29$\,minutes wall-clock per evaluation
on this hardware).
We verified that the 30-chunk estimate produces stable \emph{rankings} across
conditions: running the baseline on the full validation set gives PPL~6.10
versus 6.275 for 30 chunks, a 2.7\% difference that does not affect any
qualitative conclusion.

% ----------------------------------------------------------------
\subsection{Clamping Implementation}
\label{app:clamping}

The weight tensor \texttt{kv\_b\_proj.weight} has shape
$\bigl[n_h(d_{qk} + d_v),\; r\bigr]$,
where rows $[0 : n_h d_{qk}]$ encode $W_K^\uparrow$ and rows
$[n_h d_{qk} : n_h(d_{qk}+d_v)]$ encode $W_V^\uparrow$.
All DeepSeek-V2/V2-Lite/V3 variants satisfy $d_{qk} = d_v = 128$; this
equality is asserted before any modification.

The intervention $W_K^\uparrow \leftarrow W_V^\uparrow$ is applied by copying
the V-block into the K-block under \texttt{torch.no\_grad()}.
For each experimental condition, original weights are saved before clamping and
restored after the PPL evaluation, so all conditions share a single model
resident in memory and only one copy of the modified weights exists at a time.

\medskip\noindent
\textbf{Layer schedules.}
The \emph{middle\_block} schedule clamps layers
$\bigl[\lfloor 0.25L \rfloor, \lfloor 0.75L \rfloor\bigr)$
(layers 6--19 for $L=27$, covering 52\% of the network).
The \emph{global} schedule clamps all $L$ layers.

\medskip\noindent
\textbf{Direction comparison.}
Both directions were evaluated to check asymmetry between the key and value
pathways (Table~\ref{tab:exp3-results}).
For the \emph{middle\_block} schedule the two directions give nearly identical
$\Delta\text{PPL}$ ($+10.10$ vs.\ $+10.72$, a 6\% difference), confirming that
K and V are interchangeable in the collapse zone.
For the \emph{global} schedule, $W_K^\uparrow \leftarrow W_V^\uparrow$
(PPL~$= 1190$) is $2.9\times$ more damaging than
$W_V^\uparrow \leftarrow W_K^\uparrow$ (PPL~$= 413$).
The asymmetry is driven by early layers (layers~0--5, high $D_\text{Gram}$)
where corrupting attention routing (replacing keys with value-derived
projections) is more harmful than corrupting value extraction.


\section{Extended Experimental Details}
\label{app:experimental-details}

\subsection{Hardware and Software}
\label{app:hardware}

All experiments were run on a single workstation with 62\,GB DDR4 RAM and an
Intel Core i7-6850K (6 cores / 12 threads, 3.6\,GHz, AVX2, no native BF16).
Four installed GTX~1080~Ti GPUs (compute capability 6.1) are incompatible with
the installed PyTorch build (2.12.0+cu130, which requires CC~$\geq$~7.5); all
computation is therefore CPU-only.
BF16 matrix operations fall back to FP32 via oneDNN.
Model weights are kept in \texttt{bfloat16} throughout
(31\,GB resident memory for DeepSeek-V2-Lite).

\medskip\noindent
\textbf{Software versions:}
Python~3.12,
PyTorch~2.12.0,
Transformers~5.3.0,
safetensors~0.7.0,
datasets (Salesforce/wikitext).
All forward passes use \texttt{OMP\_NUM\_THREADS=12} and \texttt{use\_cache=False}.

% ----------------------------------------------------------------
\subsection{Model Loading and the \texttt{\_tied\_weights\_keys} Bug}
\label{app:loading-bug}

Loading DeepSeek-V2-Lite via \texttt{AutoModelForCausalLM.from\_pretrained}
under Transformers~5.x silently fails to load any weights from the checkpoint.
The root cause is in the model's remote code (\texttt{modeling\_deepseek.py}):

\begin{verbatim}
_tied_weights_keys = ["lm_head.weight"]   # old Transformers 4.x list format
\end{verbatim}

Transformers~5.x calls \texttt{.keys()} on this attribute, expecting a
\texttt{dict}; the \texttt{list} raises \texttt{AttributeError}.
The exception is caught internally and stored as a \texttt{SkipParameters}
conversion error, after which \texttt{\_initialize\_missing\_keys} calls
\texttt{\_init\_weights}, re-initialising every \texttt{nn.Linear} and
\texttt{nn.Embedding} weight to $\mathcal{N}(0,\,0.02)$.
\texttt{nn.LayerNorm} weights are not handled by \texttt{\_init\_weights}
and remain at the PyTorch default (all-ones, $\sigma=0$).
The result is that all 5{,}291 named parameters take initialisation values
rather than the trained checkpoint.

\medskip\noindent
\textbf{Symptom.}
WikiText-2 perplexity $\approx 129{,}000$.
A standard detection heuristic---flagging parameters whose empirical standard
deviation satisfies $|\sigma - 0.02| < 0.003$---corrected the
\texttt{nn.Linear} and \texttt{nn.Embedding} parameters but missed the 82
LayerNorm weights, whose trained standard deviations ($\approx 0.017$--$0.019$)
lie within the same tolerance window as the initialisation value.
After heuristic repair, perplexity remained $\approx 1.7 \times 10^{11}$,
indicating the LayerNorm weights are load-bearing.

\medskip\noindent
\textbf{Fix.}
After \texttt{from\_pretrained}, every named parameter is unconditionally
overwritten by re-reading directly from the safetensors shards:

\begin{verbatim}
for shard_file, keys in sorted(shard_to_keys.items()):
    with safe_open(shard_path, framework="pt", device="cpu") as f:
        for key in keys:
            model_params[key].data.copy_(
                f.get_tensor(key).to(dtype))
\end{verbatim}

\noindent
Shards are processed sequentially; peak extra memory is one open shard
($\approx$8\,GB) alongside the resident model ($\approx$31\,GB), well within
the 62\,GB budget.
Correctness was verified by comparing per-parameter standard deviations
against checkpoint values for a representative set spanning all weight types
(embeddings, attention projections, MoE experts, layer norms, LM head).
After the fix, WikiText-2 perplexity is 6.275 (Table~\ref{tab:exp3-results}).

% ----------------------------------------------------------------
\subsection{Perplexity Evaluation Protocol}
\label{app:ppl}

\textbf{Dataset.}
WikiText-2 validation split~\citep{merity2017pointer}, loaded via the
Salesforce/wikitext HuggingFace repository.
The full validation text is concatenated into a single token sequence
(${\approx}273$K tokens under the DeepSeek-V2-Lite tokenizer).

\medskip\noindent
\textbf{Sliding-window protocol.}
Context window $L = 512$, stride $s = 512$ (non-overlapping windows).
Each forward pass scores the final $\min(s, \text{chunk length})$ tokens.
No overlap is used; the estimate is therefore a strict upper bound on the
true held-out NLL but is consistent across all experimental conditions.

\medskip\noindent
\textbf{Sample budget.}
For Experiment~2 (activation alignment), forward passes cover the first
50{,}000 tokens (${\approx}97$ chunks).
For Experiment~3 (clamping), each PPL estimate uses the first 30 chunks
(${\approx}15{,}360$ tokens, ${\approx}29$\,minutes wall-clock per evaluation
on this hardware).
We verified that the 30-chunk estimate produces stable \emph{rankings} across
conditions: running the baseline on the full validation set gives PPL~6.10
versus 6.275 for 30 chunks, a 2.7\% difference that does not affect any
qualitative conclusion.

% ----------------------------------------------------------------
\subsection{Clamping Implementation}
\label{app:clamping}

The weight tensor \texttt{kv\_b\_proj.weight} has shape
$\bigl[n_h(d_{qk} + d_v),\; r\bigr]$,
where rows $[0 : n_h d_{qk}]$ encode $W_K^\uparrow$ and rows
$[n_h d_{qk} : n_h(d_{qk}+d_v)]$ encode $W_V^\uparrow$.
All DeepSeek-V2/V2-Lite/V3 variants satisfy $d_{qk} = d_v = 128$; this
equality is asserted before any modification.

The intervention $W_K^\uparrow \leftarrow W_V^\uparrow$ is applied by copying
the V-block into the K-block under \texttt{torch.no\_grad()}.
For each experimental condition, original weights are saved before clamping and
restored after the PPL evaluation, so all conditions share a single model
resident in memory and only one copy of the modified weights exists at a time.

\medskip\noindent
\textbf{Layer schedules.}
The \emph{middle\_block} schedule clamps layers
$\bigl[\lfloor 0.25L \rfloor, \lfloor 0.75L \rfloor\bigr)$
(layers 6--19 for $L=27$, covering 52\% of the network).
The \emph{global} schedule clamps all $L$ layers.

\medskip\noindent
\textbf{Direction comparison.}
Both directions were evaluated to check asymmetry between the key and value
pathways (Table~\ref{tab:exp3-results}).
For the \emph{middle\_block} schedule the two directions give nearly identical
$\Delta\text{PPL}$ ($+10.10$ vs.\ $+10.72$, a 6\% difference), confirming that
K and V are interchangeable in the collapse zone.
For the \emph{global} schedule, $W_K^\uparrow \leftarrow W_V^\uparrow$
(PPL~$= 1190$) is $2.9\times$ more damaging than
$W_V^\uparrow \leftarrow W_K^\uparrow$ (PPL~$= 413$).
The asymmetry is driven by early layers (layers~0--5, high $D_\text{Gram}$)
where corrupting attention routing (replacing keys with value-derived
projections) is more harmful than corrupting value extraction.


\section{Extended Experimental Details}
\label{app:experimental-details}

\subsection{Hardware and Software}
\label{app:hardware}

All experiments were run on a single workstation with 62\,GB DDR4 RAM and an
Intel Core i7-6850K (6 cores / 12 threads, 3.6\,GHz, AVX2, no native BF16).
Four installed GTX~1080~Ti GPUs (compute capability 6.1) are incompatible with
the installed PyTorch build (2.12.0+cu130, which requires CC~$\geq$~7.5); all
computation is therefore CPU-only.
BF16 matrix operations fall back to FP32 via oneDNN.
Model weights are kept in \texttt{bfloat16} throughout
(31\,GB resident memory for DeepSeek-V2-Lite).

\medskip\noindent
\textbf{Software versions:}
Python~3.12,
PyTorch~2.12.0,
Transformers~5.3.0,
safetensors~0.7.0,
datasets (Salesforce/wikitext).
All forward passes use \texttt{OMP\_NUM\_THREADS=12} and \texttt{use\_cache=False}.

% ----------------------------------------------------------------
\subsection{Model Loading and the \texttt{\_tied\_weights\_keys} Bug}
\label{app:loading-bug}

Loading DeepSeek-V2-Lite via \texttt{AutoModelForCausalLM.from\_pretrained}
under Transformers~5.x silently fails to load any weights from the checkpoint.
The root cause is in the model's remote code (\texttt{modeling\_deepseek.py}):

\begin{verbatim}
_tied_weights_keys = ["lm_head.weight"]   # old Transformers 4.x list format
\end{verbatim}

Transformers~5.x calls \texttt{.keys()} on this attribute, expecting a
\texttt{dict}; the \texttt{list} raises \texttt{AttributeError}.
The exception is caught internally and stored as a \texttt{SkipParameters}
conversion error, after which \texttt{\_initialize\_missing\_keys} calls
\texttt{\_init\_weights}, re-initialising every \texttt{nn.Linear} and
\texttt{nn.Embedding} weight to $\mathcal{N}(0,\,0.02)$.
\texttt{nn.LayerNorm} weights are not handled by \texttt{\_init\_weights}
and remain at the PyTorch default (all-ones, $\sigma=0$).
The result is that all 5{,}291 named parameters take initialisation values
rather than the trained checkpoint.

\medskip\noindent
\textbf{Symptom.}
WikiText-2 perplexity $\approx 129{,}000$.
A standard detection heuristic---flagging parameters whose empirical standard
deviation satisfies $|\sigma - 0.02| < 0.003$---corrected the
\texttt{nn.Linear} and \texttt{nn.Embedding} parameters but missed the 82
LayerNorm weights, whose trained standard deviations ($\approx 0.017$--$0.019$)
lie within the same tolerance window as the initialisation value.
After heuristic repair, perplexity remained $\approx 1.7 \times 10^{11}$,
indicating the LayerNorm weights are load-bearing.

\medskip\noindent
\textbf{Fix.}
After \texttt{from\_pretrained}, every named parameter is unconditionally
overwritten by re-reading directly from the safetensors shards:

\begin{verbatim}
for shard_file, keys in sorted(shard_to_keys.items()):
    with safe_open(shard_path, framework="pt", device="cpu") as f:
        for key in keys:
            model_params[key].data.copy_(
                f.get_tensor(key).to(dtype))
\end{verbatim}

\noindent
Shards are processed sequentially; peak extra memory is one open shard
($\approx$8\,GB) alongside the resident model ($\approx$31\,GB), well within
the 62\,GB budget.
Correctness was verified by comparing per-parameter standard deviations
against checkpoint values for a representative set spanning all weight types
(embeddings, attention projections, MoE experts, layer norms, LM head).
After the fix, WikiText-2 perplexity is 6.275 (Table~\ref{tab:exp3-results}).

% ----------------------------------------------------------------
\subsection{Perplexity Evaluation Protocol}
\label{app:ppl}

\textbf{Dataset.}
WikiText-2 validation split~\citep{merity2017pointer}, loaded via the
Salesforce/wikitext HuggingFace repository.
The full validation text is concatenated into a single token sequence
(${\approx}273$K tokens under the DeepSeek-V2-Lite tokenizer).

\medskip\noindent
\textbf{Sliding-window protocol.}
Context window $L = 512$, stride $s = 512$ (non-overlapping windows).
Each forward pass scores the final $\min(s, \text{chunk length})$ tokens.
No overlap is used; the estimate is therefore a strict upper bound on the
true held-out NLL but is consistent across all experimental conditions.

\medskip\noindent
\textbf{Sample budget.}
For Experiment~2 (activation alignment), forward passes cover the first
50{,}000 tokens (${\approx}97$ chunks).
For Experiment~3 (clamping), each PPL estimate uses the first 30 chunks
(${\approx}15{,}360$ tokens, ${\approx}29$\,minutes wall-clock per evaluation
on this hardware).
We verified that the 30-chunk estimate produces stable \emph{rankings} across
conditions: running the baseline on the full validation set gives PPL~6.10
versus 6.275 for 30 chunks, a 2.7\% difference that does not affect any
qualitative conclusion.

% ----------------------------------------------------------------
\subsection{Clamping Implementation}
\label{app:clamping}

The weight tensor \texttt{kv\_b\_proj.weight} has shape
$\bigl[n_h(d_{qk} + d_v),\; r\bigr]$,
where rows $[0 : n_h d_{qk}]$ encode $W_K^\uparrow$ and rows
$[n_h d_{qk} : n_h(d_{qk}+d_v)]$ encode $W_V^\uparrow$.
All DeepSeek-V2/V2-Lite/V3 variants satisfy $d_{qk} = d_v = 128$; this
equality is asserted before any modification.

The intervention $W_K^\uparrow \leftarrow W_V^\uparrow$ is applied by copying
the V-block into the K-block under \texttt{torch.no\_grad()}.
For each experimental condition, original weights are saved before clamping and
restored after the PPL evaluation, so all conditions share a single model
resident in memory and only one copy of the modified weights exists at a time.

\medskip\noindent
\textbf{Layer schedules.}
The \emph{middle\_block} schedule clamps layers
$\bigl[\lfloor 0.25L \rfloor, \lfloor 0.75L \rfloor\bigr)$
(layers 6--19 for $L=27$, covering 52\% of the network).
The \emph{global} schedule clamps all $L$ layers.

\medskip\noindent
\textbf{Direction comparison.}
Both directions were evaluated to check asymmetry between the key and value
pathways (Table~\ref{tab:exp3-results}).
For the \emph{middle\_block} schedule the two directions give nearly identical
$\Delta\text{PPL}$ ($+10.10$ vs.\ $+10.72$, a 6\% difference), confirming that
K and V are interchangeable in the collapse zone.
For the \emph{global} schedule, $W_K^\uparrow \leftarrow W_V^\uparrow$
(PPL~$= 1190$) is $2.9\times$ more damaging than
$W_V^\uparrow \leftarrow W_K^\uparrow$ (PPL~$= 413$).
The asymmetry is driven by early layers (layers~0--5, high $D_\text{Gram}$)
where corrupting attention routing (replacing keys with value-derived
projections) is more harmful than corrupting value extraction.


\section{Extended Experimental Details}
\label{app:experimental-details}

\subsection{Hardware and Software}
\label{app:hardware}

All experiments were run on a single workstation with 62\,GB DDR4 RAM and an
Intel Core i7-6850K (6 cores / 12 threads, 3.6\,GHz, AVX2, no native BF16).
Four installed GTX~1080~Ti GPUs (compute capability 6.1) are incompatible with
the installed PyTorch build (2.12.0+cu130, which requires CC~$\geq$~7.5); all
computation is therefore CPU-only.
BF16 matrix operations fall back to FP32 via oneDNN.
Model weights are kept in \texttt{bfloat16} throughout
(31\,GB resident memory for DeepSeek-V2-Lite).

\medskip\noindent
\textbf{Software versions:}
Python~3.12,
PyTorch~2.12.0,
Transformers~5.3.0,
safetensors~0.7.0,
datasets (Salesforce/wikitext).
All forward passes use \texttt{OMP\_NUM\_THREADS=12} and \texttt{use\_cache=False}.

% ----------------------------------------------------------------
\subsection{Model Loading and the \texttt{\_tied\_weights\_keys} Bug}
\label{app:loading-bug}

Loading DeepSeek-V2-Lite via \texttt{AutoModelForCausalLM.from\_pretrained}
under Transformers~5.x silently fails to load any weights from the checkpoint.
The root cause is in the model's remote code (\texttt{modeling\_deepseek.py}):

\begin{verbatim}
_tied_weights_keys = ["lm_head.weight"]   # old Transformers 4.x list format
\end{verbatim}

Transformers~5.x calls \texttt{.keys()} on this attribute, expecting a
\texttt{dict}; the \texttt{list} raises \texttt{AttributeError}.
The exception is caught internally and stored as a \texttt{SkipParameters}
conversion error, after which \texttt{\_initialize\_missing\_keys} calls
\texttt{\_init\_weights}, re-initialising every \texttt{nn.Linear} and
\texttt{nn.Embedding} weight to $\mathcal{N}(0,\,0.02)$.
\texttt{nn.LayerNorm} weights are not handled by \texttt{\_init\_weights}
and remain at the PyTorch default (all-ones, $\sigma=0$).
The result is that all 5{,}291 named parameters take initialisation values
rather than the trained checkpoint.

\medskip\noindent
\textbf{Symptom.}
WikiText-2 perplexity $\approx 129{,}000$.
A standard detection heuristic---flagging parameters whose empirical standard
deviation satisfies $|\sigma - 0.02| < 0.003$---corrected the
\texttt{nn.Linear} and \texttt{nn.Embedding} parameters but missed the 82
LayerNorm weights, whose trained standard deviations ($\approx 0.017$--$0.019$)
lie within the same tolerance window as the initialisation value.
After heuristic repair, perplexity remained $\approx 1.7 \times 10^{11}$,
indicating the LayerNorm weights are load-bearing.

\medskip\noindent
\textbf{Fix.}
After \texttt{from\_pretrained}, every named parameter is unconditionally
overwritten by re-reading directly from the safetensors shards:

\begin{verbatim}
for shard_file, keys in sorted(shard_to_keys.items()):
    with safe_open(shard_path, framework="pt", device="cpu") as f:
        for key in keys:
            model_params[key].data.copy_(
                f.get_tensor(key).to(dtype))
\end{verbatim}

\noindent
Shards are processed sequentially; peak extra memory is one open shard
($\approx$8\,GB) alongside the resident model ($\approx$31\,GB), well within
the 62\,GB budget.
Correctness was verified by comparing per-parameter standard deviations
against checkpoint values for a representative set spanning all weight types
(embeddings, attention projections, MoE experts, layer norms, LM head).
After the fix, WikiText-2 perplexity is 6.275 (Table~\ref{tab:exp3-results}).

% ----------------------------------------------------------------
\subsection{Perplexity Evaluation Protocol}
\label{app:ppl}

\textbf{Dataset.}
WikiText-2 validation split~\citep{merity2017pointer}, loaded via the
Salesforce/wikitext HuggingFace repository.
The full validation text is concatenated into a single token sequence
(${\approx}273$K tokens under the DeepSeek-V2-Lite tokenizer).

\medskip\noindent
\textbf{Sliding-window protocol.}
Context window $L = 512$, stride $s = 512$ (non-overlapping windows).
Each forward pass scores the final $\min(s, \text{chunk length})$ tokens.
No overlap is used; the estimate is therefore a strict upper bound on the
true held-out NLL but is consistent across all experimental conditions.

\medskip\noindent
\textbf{Sample budget.}
For Experiment~2 (activation alignment), forward passes cover the first
50{,}000 tokens (${\approx}97$ chunks).
For Experiment~3 (clamping), each PPL estimate uses the first 30 chunks
(${\approx}15{,}360$ tokens, ${\approx}29$\,minutes wall-clock per evaluation
on this hardware).
We verified that the 30-chunk estimate produces stable \emph{rankings} across
conditions: running the baseline on the full validation set gives PPL~6.10
versus 6.275 for 30 chunks, a 2.7\% difference that does not affect any
qualitative conclusion.

% ----------------------------------------------------------------
\subsection{Clamping Implementation}
\label{app:clamping}

The weight tensor \texttt{kv\_b\_proj.weight} has shape
$\bigl[n_h(d_{qk} + d_v),\; r\bigr]$,
where rows $[0 : n_h d_{qk}]$ encode $W_K^\uparrow$ and rows
$[n_h d_{qk} : n_h(d_{qk}+d_v)]$ encode $W_V^\uparrow$.
All DeepSeek-V2/V2-Lite/V3 variants satisfy $d_{qk} = d_v = 128$; this
equality is asserted before any modification.

The intervention $W_K^\uparrow \leftarrow W_V^\uparrow$ is applied by copying
the V-block into the K-block under \texttt{torch.no\_grad()}.
For each experimental condition, original weights are saved before clamping and
restored after the PPL evaluation, so all conditions share a single model
resident in memory and only one copy of the modified weights exists at a time.

\medskip\noindent
\textbf{Layer schedules.}
The \emph{middle\_block} schedule clamps layers
$\bigl[\lfloor 0.25L \rfloor, \lfloor 0.75L \rfloor\bigr)$
(layers 6--19 for $L=27$, covering 52\% of the network).
The \emph{global} schedule clamps all $L$ layers.

\medskip\noindent
\textbf{Direction comparison.}
Both directions were evaluated to check asymmetry between the key and value
pathways (Table~\ref{tab:exp3-results}).
For the \emph{middle\_block} schedule the two directions give nearly identical
$\Delta\text{PPL}$ ($+10.10$ vs.\ $+10.72$, a 6\% difference), confirming that
K and V are interchangeable in the collapse zone.
For the \emph{global} schedule, $W_K^\uparrow \leftarrow W_V^\uparrow$
(PPL~$= 1190$) is $2.9\times$ more damaging than
$W_V^\uparrow \leftarrow W_K^\uparrow$ (PPL~$= 413$).
The asymmetry is driven by early layers (layers~0--5, high $D_\text{Gram}$)
where corrupting attention routing (replacing keys with value-derived
projections) is more harmful than corrupting value extraction.
